\section{Introduction}\label{sec:intro}

Large language model (LLM) agents are increasingly used to propose supervisory actions for engineered systems:
recovery setpoints after a process fault, operating sequences, or plans that a lower-level controller then executes
\citep{yao2023react,shinn2023reflexion,car_todo}. Because an LLM's proposal can be wrong in ways that are hard to
anticipate, such agents are usually deployed behind a \emph{validator}: every proposal is checked by a simulator or
digital twin before it reaches the plant, and rejected proposals are returned to the agent with feedback. This
propose--validate--retry loop makes the agent safe by construction, but it makes every decision pay for a full
validation. Validation is often the expensive step: a digital-twin rollout must simulate the plant over the recovery
horizon, and its cost grows with the horizon, the model fidelity and the number of scenarios checked, while the cost
of the agent's own forward pass is fixed for a given prompt.

This raises a practical question: \emph{when can the validator be skipped, and how do we know?} A screen that
predicts, before validation, whether a proposal will pass could execute confident proposals directly, send clearly
bad proposals straight back to the agent, and validate only the uncertain ones. For such a screen to be useful in an
engineering setting, three things must hold. Its accept decisions must come with a guarantee on the failure rate of
what is executed unchecked; its decisions must remain safe when used inside the closed loop, not only on a static test
set; and its own cost must be small relative to the validation it replaces.

Two families of signals are available to such a screen. \emph{Observable} features describe the task state and the
proposed action, for example the plant measurements and the size of the proposed setpoint change; they are cheap but
domain specific. \emph{Model-internal} signals come from the agent itself: the confidence of its output tokens
\citep{kadavath2022know}, its hidden states \citep{azaria2023internal,burns2023ccs}, and its attention. Internal
signals are attractive because they are available for any task without a model of the plant, and because attention
offers a mechanistic reading: if the model did not attend to the part of the prompt that decides the answer, its answer
is suspect, a principle recently exploited for detecting contextual hallucinations \citep{chuang2024lookback}. We
call this signal \emph{region-based attention grounding}.

This paper studies, in a pre-registered and certification-oriented design, how many validator calls each signal lets
an LLM agent skip, at what risk, and at what cost. We use two domains with a validator in the loop: a graph
path-planning task in which every proposal is checked exactly, and fault recovery for a continuous stirred-tank
reactor (CSTR) in which proposals are validated by a digital twin \citep{car_todo}. Five open-weight models from three
families (Qwen2.5 1.5B, 3B and 7B \citep{qwen2025}, Llama-3.2-3B \citep{llama3} and SmolLM2-1.7B
\citep{allal2025smollm2}) generate the proposals.

Our contributions are as follows.
\begin{enumerate}
\item \textbf{A skip-the-validator protocol with certified accepts.} Each proposal is accepted unchecked, rejected
unchecked, or validated, using thresholds frozen on development data and a one-sided Clopper--Pearson bound on the
failure rate among unchecked accepts (Section~\ref{sec:method}). All hypotheses, rules and amendments were
pre-registered, and every test partition was evaluated once.
\item \textbf{A comparison of signals across two domains and five models.} Learned screens skip 16--91\% of validator
calls. Model internals add 5--15 points for four of five models in path planning, but only for one to three models,
each in a different way, in CSTR fault recovery (Section~\ref{sec:results}).
\item \textbf{A mechanistic account of when internals help.} Low attention on the deciding prompt region predicts
failure strongly in path planning (AUROC 0.75--0.88), where errors come from not consulting the relevant information,
and weakly or not at all in CSTR (0.48--0.72), where the model attends to the relevant physics and reasons
incorrectly.
\item \textbf{Closed-loop evidence and a deployment rule.} On 400 fresh CSTR episodes per model, trained screens
executed far fewer failing actions than random skipping at matched rates. Probes that use internal signals failed on
retries, whose prompts contain failure feedback, for two models; allowing unchecked execution only for first
proposals removed these failures with no measurable loss of recovery (Section~\ref{sec:closedloop}).
\item \textbf{A cost analysis that ties savings to simulator fidelity.} Because decisions do not depend on validator
latency, the logged closed-loop decisions can be re-costed exactly. Savings are small or negative for a fast
validator and grow to 21--64\% of total compute at 60\,s per validation for four of five models; we report the
break-even validator cost for every screen.
\end{enumerate}

The remainder of the paper reviews related work (Section~\ref{sec:related}), defines the screens, signals and
certification procedure (Section~\ref{sec:method}), describes the domains and data (Section~\ref{sec:setup}), and
reports offline results (Section~\ref{sec:results}), closed-loop results and costs (Section~\ref{sec:closedloop}).
Section~\ref{sec:discussion} discusses implications and limitations.
